% Ad Familiares 5.17 (ad-familiares-05-17)
% Source: https://raw.githubusercontent.com/PerseusDL/canonical-latinLit/master/data/phi0474/phi056/phi0474.phi056.perseus-lat1.xml
% Fetched by scripts/fetch_latin.py

% Dateline (Perseus): Scr. Romae a. 697 (57) .

\ciceroLetterOpener{M. CICERO S. D. P. SITTIO P. F.}

\ciceroSection{1} non oblivione amicitiae nostrae neque intermissione consuetudinis meae superioribus temporibus ad te nullas litteras misi, sed quod priora tempora in ruinis rei p. nostrisque iacuerunt, posteriora autem me a scribendo tuis iniustissimis atque acerbissimis incommodis retardarunt. Cum vero et intervallum iam satis longum fuisset et tuam virtutem animique magnitudinem diligentius essem mecum recordatus, non putavi esse alienum institutis meis haec ad te scribere.

\ciceroSection{2} ego te, P. Sitti, et primis temporibus illis, quibus in invidiam absens et in crimen vocabare, defendi et, cum in tui familiarissimi iudicio ac periculo tuum crimen coniungeretur, ut potui accuratissime te tuamque causam tutatus sum et proxime recenti adventu meo, cum rem aliter institutam offendissem ac mihi placuisset, si adfuissem, tamen nulla re saluti tuae defui; cumque eo tempore invidia annonae, inimici non solum tui, verum etiam amicorum tuorum, iniquitas totius iudici multaque alia rei p. vitia plus quam causa ipsa veritasque valuissent, Publio tuo neque opera neque consilio neque labore neque gratia neque testimonio defui.

\ciceroSection{3} quam ob rem omnibus officiis amicitiae diligenter a me sancteque servatis ne hoc quidem praetermittendum esse duxi, te ut hortarer rogaremque, ut et hominem te et virum esse meminisses, id est ut et communem incertumque casum, quem neque vitare quisquam nostrum nec praestare ullo pacto potest, sapienter ferres et dolori fortiter ac fortunae resisteres cogitaresque et in nostra civitate et in ceteris, quae rerum potitae sunt multis fortissimis atque optimis viris iniustis iudiciis talis casus incidisse. illud utinam ne vere scriberem, ea te re p. carere, in qua neminem prudentem hominem res ulla delectet !

\ciceroSection{4} de tuo autem filio, vereor ne, si nihil ad te scripserim, debitum eius virtuti videar testimonium non dedisse, sin autem omnia, quae sentio, perscripserim, ne refricem meis litteris desiderium ac dolorem tuum. sed tamen prudentissime facies, si illius pietatem, virtutem, industriam, ubicumque eris, tuam esse, tecum esse duces; nec enim minus nostra sunt quae animo complectimur quam quae oculis intuemur.

\ciceroSection{5} quam ob rem et illius eximia virtus summusque in te amor magnae tibi consolationi debet esse et nos ceterique, qui te non ex fortuna, sed ex virtute tua pendimus semperque pendemus, et maxime animi tui conscientia, cum tibi nihil merito accidisse reputabis et illud adiunges, homines sapientis turpitudine, non casu et delicto suo, non aliorum iniuria commoveri. ego et memoria nostrae veteris amicitiae et virtute atque observantia fili tui monitus nullo loco dero neque ad consolandam neque ad levandam fortunam tuam. tu si quid ad me forte scripseris, perficiam ne te frustra scripsisse arbitrere.
